\documentclass[../../main.tex]{subfiles}
\begin{document}

% SOURCE p=224; printed=209; object=ch06-notes-duality-history; region=middle
\textbf{关于\cref{section:ch06-lp-dual}。}Hewitt 与 Stromberg\cite[第15节]{folland-bib-076}采用了一条很不相同的途径来建立 $1<p<\infty$ 时的 $L^p$ 对偶理论。J.~Schwartz\cite{folland-bib-130}给出了适用于任意测度空间的 $(L^1)^*$ 刻画。

% SOURCE p=224; printed=209; object=ch06-notes-linfty-ba-representation; region=middle-lower
\cref{theorem:ch06-lp-duality}的证明在 $p=\infty$ 时失效，因为集合函数 $\nu(E)=\phi(\chi_E)$ 不一定可数可加。不过，它仍是 $(X,\mathscr M)$ 上有界的有限可加复测度，并且在下述意义下关于 $\mu$ 绝对连续：只要 $\mu(E)=0$，就有 $\nu(E)=0$。反过来，给定 $(X,\mathscr M)$ 上有界的有限可加复测度 $\nu$，可以定义有界可测函数关于 $\nu$ 的积分：先对简单函数以自然方式定义，再证明 $|\int f\,\mathrm d\nu|\leqslant C\|f\|_u$，从而把积分延拓到简单函数的所有一致极限。这样便把 $(L^\infty)^*$ 表示为一个有限可加复测度空间。参见 Hewitt 与 Stromberg\cite[第20节]{folland-bib-076}；有限可加积分更一般的讨论见 Dunford 与 Schwartz\cite[第3章]{folland-bib-035}。
% SUPPLEMENT: id=v14-ch06-ba-definition; parent=ch06-notes-linfty-ba-representation; kind=definition+proof+motivation+hypotheses; origin=independent-derivation
\begin{definition}[有限可加复测度及其变差（重构新增）]\label{definition:ch06-ba-measures}
设 $\mathscr M$ 为 $X$ 上的 $\sigma$ 代数。若 $\nu:\mathscr M\to\mathbb C$ 满足 $\nu(\varnothing)=0$，并对不交的 $E,F\in\mathscr M$ 满足 $\nu(E\cup F)=\nu(E)+\nu(F)$，称 $\nu$ 为\textbf{有限可加复测度}。定义其全变差范数为
\begin{align*}
\|\nu\|_{\mathrm{ba}}
=\sup\left\{\sum_{j=1}^N|\nu(E_j)|:
X=\bigsqcup_{j=1}^NE_j,\ E_j\in\mathscr M,\ N<\infty\right\}.
\end{align*}
记 $\mathrm{ba}_\mu(\mathscr M)$ 为全变差有限、且在每个 $\mu$ 零集上取零的此类 $\nu$ 所成空间。
\end{definition}
\begin{remark}
\textbf{有界性的等价含义与动机（重构新增）。}原文的“有界”也可以用 $M=\sup_{E\in\mathscr M}|\nu(E)|<\infty$ 表述，两者等价。显然 $M\leqslant\|\nu\|_{\mathrm{ba}}$。反过来，对一个有限分割，按 $\operatorname{Re}\nu(E_j)$ 的正负将分块合并，有限可加性给出
\begin{align*}
\sum_j|\operatorname{Re}\nu(E_j)|\leqslant2M.
\end{align*}
对虚部同样有此界，故 $\sum_j|\nu(E_j)|\leqslant4M$，从而全变差有限。这也说明不能只要求每个 $\nu(E)$ 单独为有限数，而遗漏对全部集合的统一界。

变差用所有有限分割测试，是为了同时控制简单函数每一块的复相位；它满足三角不等式和绝对齐次性，而范数为零时每个 $\nu(E)$ 都为零，故确为范数。只要求有限可加性，正是为了容纳 $L^\infty$ 范数不能检测的递减集合极限；在 $\mu$ 零集上消失，则是积分只依赖几乎处处等价类的必要条件。这里的“绝对连续”仅指这项零集条件，不能未经额外假设就套用可数可加测度的 Radon–Nikodym 结论。
\end{remark}
% END SUPPLEMENT: id=v14-ch06-ba-definition
% SUPPLEMENT: id=v14-ch06-ba-isometry; parent=ch06-notes-linfty-ba-representation; kind=proof+motivation+hypotheses; origin=independent-derivation
\begin{theorem}[$L^\infty$ 对偶的有限可加表示（重构新增）]\label{theorem:ch06-linfty-ba-isometry}
对任意测度空间 $(X,\mathscr M,\mu)$，映射
\begin{align*}
(L^\infty(\mu))^*&\longrightarrow\mathrm{ba}_\mu(\mathscr M),&
\phi&\longmapsto\nu_\phi,\qquad\nu_\phi(E)=\phi(\chi_E)
\end{align*}
是等距线性同构。其逆映射为 $\nu\mapsto(f\mapsto\int f\,\mathrm d\nu)$，积分以简单函数的一致逼近定义。
\end{theorem}
\begin{proof}
给定 $\phi$，示性函数在不交并上的加法公式及 $\phi$ 的线性性给出 $\nu_\phi$ 的有限可加性。若 $\mu(E)=0$，则 $\chi_E$ 为 $L^\infty$ 中的零元，故 $\nu_\phi(E)=0$。对有限分割 $(E_j)$，在 $\nu_\phi(E_j)\ne0$ 时取 $c_j=\overline{\nu_\phi(E_j)}/|\nu_\phi(E_j)|$，否则取 $c_j=0$。于是 $s=\sum_jc_j\chi_{E_j}$ 满足 $\|s\|_\infty\leqslant1$，并且
\begin{align*}
\sum_j|\nu_\phi(E_j)|=\phi(s)\leqslant\|\phi\|.
\end{align*}
左式为非负实数，所以最后不等式来自泛函的绝对值界。取上确界得 $\|\nu_\phi\|_{\mathrm{ba}}\leqslant\|\phi\|$。

反过来，设 $\nu\in\mathrm{ba}_\mu(\mathscr M)$。对按不交可测分割写出的简单函数 $s=\sum_jc_j\chi_{E_j}$，定义 $I_\nu(s)=\sum_jc_j\nu(E_j)$，必要时添加系数为零的剩余块。任意两种表示可用交集分割共同细化，有限可加性保证所得值相同，并保证线性性。而
\begin{align*}
|I_\nu(s)|\leqslant\sum_j|c_j|\,|\nu(E_j)|
\leqslant\|s\|_u\|\nu\|_{\mathrm{ba}}.
\end{align*}
每个有界可测复函数 $f$ 都可由简单函数一致逼近：分别把其实部、虚部的有界值域切成网格宽度趋零的有限区间，并在各可测原像上取相应网格值。若 $s_n\to f$ 一致，上式使 $I_\nu(s_n)$ 为 Cauchy 列；定义 $I_\nu(f)$ 为其极限。两组逼近之差一致趋零，故极限与选取无关。极限过程还给出线性性和 $|I_\nu(f)|\leqslant\|f\|_u\|\nu\|_{\mathrm{ba}}$。

若有界可测 $h$ 在某个 $\mu$ 零集 $N$ 之外为零，则上述简单逼近也可在 $N$ 之外取零，其所有非零分块均为 $\mu$ 零集，故积分为零；于是 $I_\nu(h)=0$。每个 $L^\infty$ 等价类可选到处有界的可测代表：把超过其本质上界的可测零集上的值改成零。不同代表的差在零集外为零，所以所得积分相同。对任意 $\varepsilon>0$ 选代表使一致范数不超过 $\|f\|_\infty+\varepsilon$，再令 $\varepsilon\downarrow0$，得到
\begin{align*}
|I_\nu(f)|\leqslant\|\nu\|_{\mathrm{ba}}\|f\|_\infty.
\end{align*}
因此 $I_\nu\in(L^\infty)^*$。由于 $I_\nu(\chi_E)=\nu(E)$，第一次构造把它送回 $\nu$。反方向上，$I_{\nu_\phi}$ 与 $\phi$ 在简单函数上相等；简单函数在 $L^\infty$ 中范数稠密，由连续性它们处处相等。结合两个方向的范数不等式，得到 $\|\nu_\phi\|_{\mathrm{ba}}=\|\phi\|$。两种构造的线性性也已由定义得到。
\end{proof}
\begin{remark}
\textbf{动机与逐条件说明（重构新增）。}每个示性函数都属于 $L^\infty$，故无需限制为有限测度集，也不需 $\sigma$ 有限性。$
\nu$ 的有限可加性确保简单积分良定义；有限变差控制一致逼近误差；零集条件确保积分下降到 $L^\infty$ 商空间。若没有最后的条件，在 $[0,1]$ 上的点质量就会把同一 $L^\infty$ 元的两个代表 $0$ 和 $\chi_{\{0\}}$ 送到不同数。测度空间不必完备，因为这里修改的集合本身已经可测。原文所说的“一致极限”先在到处有界的可测代表上进行，再由零集条件转为本质上界范数；两步不能混淆。
\end{remark}
% END SUPPLEMENT: id=v14-ch06-ba-isometry

% SOURCE p=224; printed=209; object=ch06-notes-point-mass-pathology; region=lower
\cref{example:ch06-linfty-point-evaluation}中给出的 $\phi\in(L^\infty)^*\setminus L^1$，表明有限可加测度可能有多么反常：若 $\nu(E)=\phi(\chi_E)$，则 $\nu\ll m$，但对任意连续函数积分时，$\nu$ 却表现得像零点处的点质量。
% SUPPLEMENT: id=v14-ch06-ba-point-mass; parent=ch06-notes-point-mass-pathology; kind=proof+motivation+hypotheses; origin=independent-derivation
\begin{remark}
\textbf{不满足可数可加性的完整核验（重构新增）。}在该例的 $[0,1]$ 上，$\phi$ 是连续函数取值泛函 $h\mapsto h(0)$ 的保范延拓。上面的表示给出 $\nu(X)=\phi(1)=1$、$\nu(\{0\})=0$，而对所有连续 $h$ 有 $\int h\,\mathrm d\nu=h(0)$。假若 $\nu$ 可数可加，它便是有限变差的通常复测度。令 $h_n(x)=\max(1-nx,0)$，则 $0\leqslant h_n\leqslant1$ 且 $h_n\to\chi_{\{0\}}$ 逐点。对有限正测度 $|\nu|$ 用控制收敛定理，得到 $\int h_n\,\mathrm d\nu\to\nu(\{0\})=0$；但每个积分都等于 $h_n(0)=1$，矛盾。这里通常复测度积分与上面的一致延拓积分在简单函数上一致，故对有界可测函数亦一致，比较合法。动机是把连续函数无法区分的“集中在零点附近”与可测单点的零测度性质放在一起；这个例子正说明零集意义的绝对连续性不能补回可数可加性。
\end{remark}
% END SUPPLEMENT: id=v14-ch06-ba-point-mass
\end{document}
